\documentclass[12pt, letterpaper]{article}

%-------Packages---------
\usepackage{amssymb,amsfonts}
\usepackage{mathptmx}
\usepackage{amsthm}
\usepackage{graphicx}
\usepackage[all,arc]{xy}
\usepackage[colorlinks=true, allcolors=blue]{hyperref}
\usepackage{enumerate}
\usepackage{bbm}
\usepackage{dsfont}
\usepackage{mathrsfs}
\usepackage{tikz-cd}
\usepackage{setspace}
\usepackage{import}
\usepackage[utf8]{inputenc}
\usepackage{hyperref}
\usepackage{tikz}
\usepackage{float}
\usepackage{mdframed}
\usepackage{fancyhdr}
\usepackage[nointegrals]{wasysym}

\usepackage[left=1in,top=1in,right=1in,bottom=1in]{geometry}

\usepackage{mathtools}
\usepackage{setspace}
\usepackage{cleveref}
\usepackage{multicol}
\usepackage{mathpazo}

%--------Theorem Environments--------
%theoremstyle{plain} --- default
\newmdtheoremenv{theorem}{Theorem}[subsection]
\newmdtheoremenv{corollary}[theorem]{Corollary}
\newtheorem{proposition}[theorem]{Proposition}
\newmdtheoremenv{lemma}[theorem]{Lemma}
\newtheorem{conj}[theorem]{Conjecture}
\newtheorem{quest}[theorem]{Question}

\theoremstyle{definition}
\newmdtheoremenv{definition}[theorem]{Definition}
\newmdtheoremenv{definitions}[theorem]{Definitions}
\newtheorem{example}[theorem]{Example}
\newtheorem{algorithm}[theorem]{Algorithm}
\newtheorem{rmk}[theorem]{Remark}
\newtheorem{exercise}[theorem]{Exercise}

%----------- my commands -------------
\newcommand{\C}{\mathbb{C}}
\newcommand{\R}{\mathbb{R}}
\newcommand{\Q}{\mathbb{Q}}
\newcommand{\Z}{\mathbb{Z}}
\newcommand{\N}{\mathbb{N}}
\renewcommand{\P}{\mathbb{P}}
\newcommand{\T}{\mathcal{T}}
\newcommand{\contr}{\Rightarrow \Leftarrow}
\newcommand{\HRule}[1]{\rule{\linewidth}{#1}}
\newcommand{\s}{\(\,\)}
\renewcommand{\t}[1]{\text{#1}}
\newcommand{\ang}[1]{\left\langle #1 \right\rangle}

% Meta data
\setstretch{1.2}

\begin{document}

\pagestyle{fancy}
\fancyhf{}
\fancyhead[LOE]{\slshape{\textbf{Vincent Ferrara}}}
\fancyhead[ROE]{\thepage}

\subsection*{Problem 1}
\begin{itemize}
    \item Let $A=\displaystyle{\bigcap_{\alpha \in J}}\T_\alpha$
    \item Since for all $\mathcal{T}_{\alpha}$ s.t. $\alpha \in J$ are topologies they all have $X \in \mathcal{T}_\alpha$ and $\emptyset \in \mathcal{T}_\alpha$ thus $X, \emptyset \in A$
    \item Let $U,V \in A$. Thus, $U,V \in \T_\alpha$ for all $\alpha \in J$ thus since these are topologies this implies that $U\cap V \in \T_\alpha$ for all $\alpha \in J$ thus $U \cap V \in A$
    \item Let $U=\displaystyle{\bigcup_{i \in I}} U_i$ s.t. $U_i \in A$ and $I$ is an indexing set. Since each $U_j \in A$ this means it is in every $\T_\alpha$. Since these are topologies they are closed under arbitrary unions thus $U \in \T_\alpha$ for all $\alpha \in J$ thus $U \in A$.
\end{itemize}
Thus, $A$ is a topology on $X$.
\paragraph*{Counter Example:}
\phantom{text}\\
Let $X=\{0,1,2\}$. Define two topologies on $X$:
\begin{itemize}
    \item $\T_1 = \{\emptyset, X, \{0\}, \{1,2\}\}$
    \item $\T_2 = \{\emptyset, X, \{1\}, \{0,2\}\}$
\end{itemize}
It is straightforward to check each is a topology.\\
Now consider
\[\T_1 \cup \T_2 = \{\emptyset, X, \{0\}, \{1,2\}, \{1\}, \{0,2\}\}\]
However:
\[\{1,2\},\{0,2\} \in \T_1 \cup \T_2\ \quad \t{but} \quad \{1,2\} \cap \{0,2\} = \{2\} \notin \T_1 \cap \T_2\]
So it is not true that unions of topologies are always topologies.

\newpage
\subsection*{Problem 2}
Assume $U_1,\dots,U_n$ is a finite collection of open subsets of a topological space $X$.
\paragraph*{Base Case:}
\phantom{text}\\
Let $n=2$, and let $U_1,U_2$ be open. Thus, $U_1 \cap U_2$ are open in $X$ by the definition of a topology.
\paragraph*{Inductive Step:}
\phantom{text}\\
Assume for all $n \leq k$ s.t. $2 \leq k$ and $n,k \in \N$ the union of $n$ open sets is open.\\
Consider $k+1$:\\
Let $U_1,\dots,U_{k+1}$ be open. Now consider $U_1 \cap \dots \cap U_k$ by the inductive hypothesis this is open call this set $U$. Thus, consider $U \cap U_{k+1}$ by the definition of topology the intersection of two open sets is open. Thus, by induction the intersection of $n$ open sets of a topological space is open for all $n \in \N$

\newpage
\subsection*{Problem 3}
\begin{itemize}
    \item Let $\{C_\alpha\}_{\alpha \in J}$ be a collection of closed sets in $X$ with indexing set $J$. Thus, by definition of closed the complment is open thus $C_\alpha=X \setminus O_\alpha$ for some open set $O_\alpha$. Consider:
    \[\bigcap_{\alpha \in J}C_\alpha=\bigcap_{\alpha \in J}(X \setminus O_\alpha)\]
    By De Morgan's Law:
    \[=X \setminus \bigcup_{\alpha \in J}O_\alpha\]
    Thus since open sets are closed under arbitrary union that is an open set thus since the complement of that intersection is open we know that the intersection is closed.
    \item \begin{itemize}
        \item Base Case: Let $n=2$. Let $C,V$ be closed. Thus $C=X\setminus O$ and $V=X \setminus U$ for open sets $O,U$. Consider: 
        \[C \cup V=(X \setminus O) \cup (X \setminus U)\]
        By De Morgan's Law:
        \[=X \setminus (O \cap U)\]
        Since $O \cap U$ is open we know that the complement of $C \cup V$ is open this $C \cup V$ is closed.
        \item Inductive Step: Assume for all unions of closed sets of less than or equal to $k\geq2$ is closed\\
        Consider $k+1$ sets. Let $C_1\dots C_{k+1}$ be closed. Thus consider $C_1 \cup \dots \cup C_k$. By the inductive hypothesis this is closed. Thus call this set $C$. Thus, consider $C \cap C_{k}$ by the base case it is also closed. Thus, unions of closed sets of size $n$ is closed for all $n \in \N$.
    \end{itemize}
\end{itemize}
Let $X=\R$ and use the standard topology on $\R$. Consider the collection $U_n = [\frac{1}{n},1]$ for $n \in \N$. These are by definition closed since it is the complement of $(-\infty,\frac{1}{n}) \cup (1,\infty)$.\\
Consider:
\[U=\bigcup_{n \in \N}U_n\]
Let $x \in (0,1]$ thus by the archimedian property there is $n \in \N$ s.t. $1/n < x$ thus $x \in U_n \subseteq U$\\
Let $x \in U$ thus $x \in U_n$ for some $n \in \N$ thus since $[1/n,1] \subset (0,1]$ thus $x \in (0,1]$ thus $U=(0,1]$\\
However since $\R \setminus U=(\infty,0] \cup (1,\infty)$ is not open. This means that $U$ is not closed. Thus, this is an example of how closed sets are not closed under arbitrary union.

\newpage
\subsection*{Problem 4}
Consider $H=\{(x,y,z) \in \R^3 \mid z >0\}$
Let $p \in H$. Let $\epsilon = z/2 > 0$\\
Consider the ball $B_\epsilon(p)$\\
Let $q \in B_\epsilon(p)$ thus $q=(x',y',z')$. To show that this is in $H$ all we need to show is that $z' >0$.\\
From the definition of a ball we know that $|z'-z| < \epsilon \implies z-\epsilon < z' \implies z-z/2=z/2 < z'$. Thus since $z>0$ we know that $z/2 > 0$ thus $0 < z'$. Thus $q \in H$.\\
Thus, for every point we can find a ball that is a subset of $H$. Thus, $H$ is an open subset of $\R^3$

\newpage
\subsection*{Problem 5}
\begin{enumerate}[(a)]
    \item \begin{itemize}
        \item Let $x \in X$ since $X$ is the unions of sets in $\mathcal{S}$ there exists an $x \in S$ s.t. $S \in \mathcal{S}$. Thus, since $S=S \cap S$, $x$ is in the union of these finite intersections.
        \item Let $U_1,U_2$ be in the set of the union of these finite intersections. Thus, $U_1 = S_1 \cap S_2 \cap \dots \cap S_n$, and $U_2 = T_1 \cap T_2 \cap \dots \cap T_m$ for some $S_i,T_j \in \mathcal{S}$ for $i \in \{1,\dots,n\}$ and $j \in \{1,\dots,m\}$\\
        Consider $U_1 \cap U_2$:\\
        Let $x \in U_1 \cap U_2$. Consider $U_1 \cap U_2 = U_3$. Since $U_1$ and $U_2$ are finite intersections this means $U_3$ is a finite intersection thus it is in our set. Thus $x \in U_3 \subseteq U_1 \cap U_2$
    \end{itemize}
    Thus this set is a basis for a topology on $X$.
    \item Let $\{\T_\alpha\}_{\alpha \in J}$ be all the topologies that contain $\mathcal{S}$. Let\\
    \[\T = \bigcap_{\alpha \in J}T_\alpha\]
    Let $B$ be the topology generated by the basis in (a).\\
    $T$ is the intersection of every topology that contains $\mathcal{S}$ and $B$ contains $\mathcal{S}$ we know that $\T \subseteq B$.\\
    Since $B$ is the collection of all unions of finite intersections of elements of $\mathcal{S}$, and from problem 1 we know it is a topology. This means we are closed under arbitrary unions and finite intersections thus every element of $B$ is in $\T$.\\
    Thus, $\T = B$\\
    Thus we know that $B$ is the coarsest topology that contains $\mathcal{S}$.
\end{enumerate}

\newpage
\subsection*{Problem 6}
\begin{enumerate}[(a)]
    \item Let $\beta_d$ be the basis for the order topology on $\R\times \R$ with dictionary order, and let $\beta_p$ be the basis for the product topology.\\
    Let $B \in \beta_p$. WLOG let $B=(a,b)\times(c,d)$ since we can create any basis element from unions or intersections of this type of element. Let $p \in B$ s.t. $p=(x,y)$. Consider the basis element $C=((x,y-\frac{y-c}{2}),(x,y+\frac{d-y}{2}))$. Thus, $p \in C \subseteq B$, and $C$ is clearly a basis element for the order topology. Thus, we have shown that the order topology is finer than the product topology.\\
    Consider the set $A=\{(x,y) \in R\times R \mid x-0 \t{ and } 0 < y < 1\}$. This is open in the order topology since it corresponds to the interval $(0,0)<(0,y)<(0,1)$. However, this is not open in the product topology, since it is equivalent to the set $A=\{0\} \times (0,1)$, and $\{0\}$ is not open in $\R$. Thus, $A$ is not open in the product topology.
    \item Consider the point $p=(x,y) \in \Z \times \Z$. Consider the open set in the order topology: $((x,y-1),(x,y+1))=\{(x,y)\}$. Thus, every singleton is actually an open set in $\Z \times \Z$. Thus, we can write every subset as a union of singletons thus we are actually in the discrete topology since every subset is open. Consider the open set in the product topology $(x-1,x+1) \times (y-1,y+1)=\{p\}$. Again we have shown all singletons are open thus all subsets are open thus we are in the discrete topology. Thus, these two topologies are the same.
    \item Consider the point $p=(x,y) \in \N \times \N$. If $x=y=0$ then consider the open set $[0,1) \times [0,1)=\{p\}$. Thus, this singleton is open. Now if $x=0$ consider the open set $(x-1,x+1) \times [0,1)=\{p\}$ now we know this singleton is open and do the same for if $x \neq 0$ and $y=0$. If they are both not zero then consider the open set $(x-1,x+1) \times (y-1,y+1)=\{p\}$. Thus, all singletons are open thus by the same argument as above we are in the discrete topology. Since the discrete topology is equal to all the subsets of $\N \times \N$. This means that any other topology will be coarser than it since any collection of subsets is a subset to the power set. Thus, the dictionary order topology is coarser than the product topology.\\
    Consider the open set $\{(1,0)\}$ in the product topology. For this to be open in the order toplogy we need a basis element s.t. $((a,b),(c,d))=\{(1,0)\}$. If $a=0$ then we will have all elements $(0,b+1),(0,b+2),\dots$. Thus, $a \neq 0$. If $c>1$, then we will have all elements of the form $(c,x)$ s.t. $x < d$. Thus, $c > 1$. Thus, $a=1=c$, since $a \leq 1 \leq c$. Thus, now consider $b,d$. If $d>1$ then we would have the element $(1,1)$ thus $d=1$. Thus we have the interval $((1,b),(1,1))$. For this set to only have $(1,0)$ this would mean that $d<0$. However that cannot happen since $d \in \N$. Thus this singleton is not open in the order topology. Thus, these two topologies are not equal.
\end{enumerate}

\newpage
\subsection*{Problem 7}
\begin{enumerate}[(a)]
    \item Let $X$ be a finite metrizable space. Let $d$ be a metric on $X$ that induces a topology on $X$.
    Let $X=\{x_1,\dots,x_n\}$. Thus define $r_i$ for each $i \in \{1,\dots,n\}$ to be the minimum value of $\{d(x_i,x_j \mid i \neq j \in \{1,\dots,n\})\}$. There exists a min since we only have a finite amount of points thus we will only have $n-1$ values.
    By the definition of metrizable we know that $r_i>0$ for all $i$. Thus consider the open set $B(x_i,r_i/2)=\{x \in X \mid d(x_i,x)<r_i/1\}$. Since $r_i/2<r_i \leq d(x_i,x_j)$ for $j \neq i$. This ball only contains the point $x_i$. Thus we have shown that all singletons are open sets. Thus we can unions any pair of singletons to get any subset of $X$. Thus, we know that all open sets are open.
    Thus, it is discrete.
    \item Look at $\Q$ with standard topology. Consider the set $\{1\}$ By definition of the standard topology on $\Q$ for $\{1\}$ to be open there must exists a $\epsilon >0$ s.t. $B(1,\epsilon)=\{x \in \Q \mid |x-1|< \epsilon\}=(1-\epsilon,1+\epsilon)=\{1\}$. However, by the archimedian property there exists a $n \in \N$ s.t. $1/n < \epsilon$. Thus consider the value $1+1/n \in \Q$. Thus since $1-\epsilon < 1+1/n < 1 + \epsilon$ this value is in our ball thus there is no
    $\epsilon >0$ s.t. $1$ is the only thing in the set. Thus, $\{1\}$ is not open. Thus, $\Q$ is not discrete under the standard topology.
    \item Let $X$ be a discrete topological space. Consider:\\
    \begin{equation*}
        d(x,y)=\begin{cases}
            0, \quad \t{if } x=y,\\
            1, \quad \t{if } x\neq y.
        \end{cases}
    \end{equation*}
    \begin{itemize}
        \item $d(x,y)=0=d(y,x) \iff x=y$ this is given from the definition of $d$.
        \item Assume $y\neq x$ and $x,y \in X$. Thus $d(x,y)=1=d(y,x)$
        \item Let $x,y,z \in X$.\\
        Assume $x \neq z$ since if this is true then $d(x,z)=0 \leq d(x,y)+d(y,z)$ is always true since $d$ is always $0$ or $1$.\\
        If $x=y\neq z$ then $d(x,z)=1 \leq d(x,y)+d(y,z)=1$\\
        If $x \neq y \neq z$ then $d(x,z)=1 \leq d(x,y)+d(y,z)=2$\\
        Thus $d$ is a metric.\\
        Thus consider the basis for the metric topology. Let $x \in X$, then consider the ball $B_\epsilon(x)$ for $\epsilon = 1/2$. Since all distances are 1 there will only be one thing in this ball and that is $x$ thus we know that all singletons are open. Thus, we can get any subset by unions of singletons. Thus, all subsets of $X$ are open. Thus $X$ is metrizable since we can induce the same topology with a metric.
    \end{itemize}
\end{enumerate}
\end{document}